% ==================================================
% ============== PŘÍKLADY K PROCVIČENÍ =============
% ==================================================

\exercisepart


% ==================================================
% =============== ALGEBRAICKÝ TVAR =================
% ==================================================

\begin{exerciseblock}[2]{Vypočítejte a výsledek zapište v algebraickém tvaru}
    \task $2+i+3(7-i)$
    \task $(1-i)(2i-3)-i$
    \task $(2+i)(3+i)(1+2i)$
    \task $(i+1)(i-1)+(2i+1)^2$
    \task $(3-2i)+(7+5i)-(4-3i)$
    \task $4(2-i)-3(1+2i)$
    \task $(2+3i)^2$
    \task $(4-i)(2+5i)$
    \task $(1-2i)^3$
    \task $(2+i)^2-(2-i)^2$
\end{exerciseblock}


\begin{exerciseblock}[2]{Vydělte komplexní čísla a výsledek zapište v algebraickém tvaru}
    \task $\displaystyle \frac{3i+1}{2+i}$
    \task $\displaystyle \frac{100}{3+4i}$
    \task $\displaystyle \frac{3i+2}{2i-3}$
    \task $\displaystyle \frac{1+4i}{i}$
    \task $\displaystyle \frac{3+i\sqrt{3}}{3-i\sqrt{3}}$
    \task $\displaystyle \frac{2-5i}{1+3i}$
    \task $\displaystyle \frac{4+i}{2-i}$
    \task $\displaystyle \frac{1-i}{1+i}$
    \task $\displaystyle \frac{7+2i}{3-2i}$
    \task $\displaystyle \frac{5i-1}{2+4i}$
\end{exerciseblock}


\begin{exerciseblock}[2]{Upravte}
    \task $\displaystyle (2+i)i+\frac{3+i}{2-i}$
    \task $\displaystyle \frac{2+i}{i}+\frac{i}{i+1}-\frac{2i+1}{i-1}$
    \task $\displaystyle \frac{i-1}{2}-\frac{i}{i-1}\cdot i+1$
    \task $\displaystyle (5i-1):\left(2-\frac{i+3}{2+i}\right)$
    \task $\displaystyle \frac{\frac{i}{2-i}+\frac1i}{1+\frac{i}{2i+1}}$
    \task $\displaystyle \frac{1+i}{1-i}+\frac{1-i}{1+i}$
    \task $\displaystyle \left(\frac{2+i}{1-i}\right)^2$
    \task $\displaystyle \frac{(1+i)^4}{(1-i)^2}$
    \task $\displaystyle \frac{3-2i}{1+i}-\frac{1+4i}{2-i}$
    \task $\displaystyle \frac{1}{1+i}+\frac{1}{1-i}+\frac{1}{i}$
\end{exerciseblock}


\begin{exercise}{Určete reálná čísla $x$ a $y$.}
    \begin{tasks}(2)
        \task $(2x+1)+(y-3)i=7+5i$
        \task $(x-y)+(x+y)i=4-2i$
        \task $(3x-2)+(2y+1)i=x+4+(y-5)i$
        \task $(x+yi)(1+i)=5+i$
        \task $\displaystyle \frac{x+yi}{1-i}=2+3i$
        \task $(x+2i)^2=5+4yi$
    \end{tasks}
\end{exercise}


\begin{exercise}{Určete hodnoty reálného parametru $k$ tak, aby dané komplexní číslo mělo požadovanou vlastnost.}
    \begin{tasks}(2)
        \task $\displaystyle z=\frac{k+2i}{k-i}$ bylo reálné,
        \task $\displaystyle z=\frac{k+2i}{k-i}$ bylo ryze imaginární,
        \task $\displaystyle z=\frac{3+ki}{1+i}$ mělo reálnou část rovnu $2$,
        \task $\displaystyle z=\frac{k-i}{2+i}$ mělo imaginární část rovnu $-1$,
        \task $(k+i)^2$ bylo reálné,
        \task $(2+ki)(1-i)$ bylo ryze imaginární.
    \end{tasks}
\end{exercise}


% ==================================================
% ============== MOCNINY JEDNOTKY i ================
% ==================================================

\subsection*{Mocniny imaginární jednotky}

\begin{exerciseblock}[5]{Vypočítejte}
    \task $i^2$
    \task $i^3$
    \task $i^4$
    \task $i^{17}$
    \task $i^{42}$
    \task $i^{50}$
    \task $i^{125}$
    \task $i^{505}$
    \task $i^{2026}$
    \task $i^{10001}$
\end{exerciseblock}


\begin{exerciseblock}[2]{Vypočítejte}
    \task $1+i+i^2+i^3+i^4$
    \task $1+i^2+i^4+i^6+i^8+i^{10}$
    \task $1+i^3+i^5+i^7+i^9+i^{11}$
    \task $1+i+i^2+\cdots+i^{20}$
    \task $i^{12}-i^{29}+5i^{6}-3i^{11}$
    \task $3i^{2025}+2i^{2026}-i^{2027}$
    \task $i^{100}+i^{101}+i^{102}+i^{103}$
    \task $2i^{38}-4i^{57}+i^{84}$
    \task $\displaystyle \sum_{k=0}^{15}i^k$
    \task $\displaystyle i+i^5+i^9+\cdots+i^{97}$
\end{exerciseblock}


\begin{exerciseblock}[2]{Vypočítejte}
    \task $i^{-1}$
    \task $i^{-2}$
    \task $i^{-3}$
    \task $i^{-37}$
    \task $i^{-78}$
    \task $i^{-1001}$
    \task $1+i^{-1}+i^{-2}+i^{-3}+i^{-4}$
    \task $i^{-30}+i^{-40}+i^{-50}+i^{-60}$
    \task $3i^{-7}-2i^{-13}$
    \task $\displaystyle \frac{i^{57}+i^{23}}{i^{14}}$
\end{exerciseblock}


\begin{exerciseblock}[2]{Vypočítejte}
    \task $(1+i)^2$
    \task $(1-i)^2$
    \task $(1+i)^3$
    \task $(1-i)^3$
    \task $(1+i)^4$
    \task $(1-i)^4$
    \task $(1+i)^8$
    \task $(1-i)^{12}$
    \task $(1+i)^{20}$
    \task $\displaystyle \frac{(1+i)^{10}}{(1-i)^6}$
\end{exerciseblock}


% ==================================================
% ================= GAUSSOVA ROVINA =================
% ==================================================

\subsection*{Gaussova rovina}

\begin{exercise}{Znázorněte v Gaussově rovině obrazy následujících komplexních čísel.}
    \[
        z_1=1+2i,\qquad
        z_2=3-i,\qquad
        z_3=-2+4i,\qquad
        z_4=-3-2i,
    \]
    \[
        z_5=5,\qquad
        z_6=-4i,\qquad
        z_7=-1+i,\qquad
        z_8=2-3i.
    \]
\end{exercise}


\begin{exercise}{Pro čísla $z_1=1+2i$ a $z_2=3-i$ nejprve určete algebraicky a poté graficky znázorněte obrazy čísel}
    \begin{tasks}(4)
        \task $z_1+z_2$
        \task $z_1-z_2$
        \task $-z_1$
        \task $2z_2$
        \task $\overline{z_1}$
        \task $\overline{z_2}$
        \task $iz_1$
        \task $-iz_2$
    \end{tasks}
\end{exercise}


\begin{exercise}{Nakreslete v Gaussově rovině množinu všech komplexních čísel $z$, pro která platí}
    \begin{tasks}(2)
        \task $|z|=3$
        \task $|z|<2$
        \task $|z|\leq4$
        \task $|z-1|=2$
        \task $|z+2i|=3$
        \task $|z-(2-i)|<2$
        \task $|z-3i|\geq2$
        \task $|z-(1+i)|\leq1$
        \task $|z-2|=|z+2|$
        \task $|z-i|=|z-3i|$
    \end{tasks}
\end{exercise}


% ==================================================
% =========== KOMPLEXNĚ SDRUŽENÁ ČÍSLA =============
% ==================================================

\subsection*{Komplexně sdružená čísla}

\begin{exerciseblock}[2]{Napište komplexně sdružená čísla}
    \task $z_1=2+i$
    \task $z_2=7-3i$
    \task $z_3=4i$
    \task $\displaystyle z_4=\frac{1+5i}{3}$
    \task $z_5=\sqrt{3}+1-7i$
    \task $z_6=\sqrt{10}+i+i\sqrt2$
    \task $z_7=-5$
    \task $z_8=-\sqrt2 i$
    \task $z_9=(2+i)(3-i)$
    \task $\displaystyle z_{10}=\frac{3+4i}{1-2i}$
\end{exerciseblock}


\begin{exerciseblock}[2]{Vypočítejte}
    \task $\overline{2+i}+1-\overline{i}$
    \task $\overline{3+4i}+3-\overline{7i}$
    \task $\overline{(1+i)(3+2i)}$
    \task $\overline{(1+i)}\cdot\overline{(3+2i)}$
    \task $\overline{(5+3i)^2}$
    \task $\left(\overline{5+3i}\right)^2$
    \task $\displaystyle \overline{\frac{1+i}{2-i}}$
    \task $\displaystyle \frac{\overline{1+i}}{\overline{2-i}}$
    \task $z+\overline z$ pro $z=3-7i$
    \task $z-\overline z$ pro $z=-2+5i$
\end{exerciseblock}


\begin{exercise}{Dokažte pro libovolná komplexní čísla $z_1,z_2$, že platí}
    \begin{tasks}(2)
        \task $\overline{\overline z}=z$
        \task $\overline{z_1+z_2}=\overline{z_1}+\overline{z_2}$
        \task $\overline{z_1z_2}=\overline{z_1}\,\overline{z_2}$
        \task $z+\overline z=2\operatorname{Re}(z)$
        \task $z-\overline z=2i\operatorname{Im}(z)$
        \task $z\overline z=|z|^2$
    \end{tasks}
\end{exercise}


% ==================================================
% ================ ABSOLUTNÍ HODNOTA ================
% ==================================================

\subsection*{Absolutní hodnota komplexního čísla}

\begin{exerciseblock}[2]{Vypočítejte}
    \task $|6+2i|$
    \task $\left|\frac{\sqrt3+1}{3}-\frac{\sqrt3-1}{3}i\right|$
    \task $|\sqrt5+2+2i-i\sqrt5|$
    \task $\left|\frac{\sqrt7}{4}(1+i)+\frac{\sqrt5}{4}(1-i)\right|$
    \task $|-3+4i|$
    \task $|5i|$
    \task $|-7|$
    \task $|1-\sqrt3 i|$
    \task $|\sqrt2+\sqrt2 i|$
    \task $\left|\frac{3-4i}{5}\right|$
\end{exerciseblock}


\begin{exerciseblock}[2]{Vypočítejte}
    \task $|(7+i)(4-3i)|$
    \task $\left|\frac{4-2i}{3+i}\right|$
    \task $\left|\frac{10i}{2\sqrt6-2\sqrt3 i}\right|$
    \task $\left|\frac{4-3i+i}{3-2i}\right|$
    \task $\left|\frac{3-4i}{5i}\right|\cdot\left|\frac{1+i}{3-i}\right|$
    \task $|(1+i)^8|$
    \task $\left|\frac{(2+i)^2}{1-2i}\right|$
    \task $|z\overline z|$ pro $z=3+4i$
    \task $|z^5|$ pro $|z|=2$
    \task $\left|\frac{z_1^3}{z_2^2}\right|$, je-li $|z_1|=2$, $|z_2|=4$
\end{exerciseblock}


\begin{exercise}{Určete všechna komplexní čísla $z$, která splňují danou podmínku, a množinu řešení znázorněte v Gaussově rovině.}
    \begin{tasks}(2)
        \task $|z|=2$
        \task $|z-1|=3$
        \task $|z+2i|\leq4$
        \task $1<|z-3+i|\leq3$
        \task $|z|=|z-4|$
        \task $|z-i|=|z+i|$
        \task $|z-1|<|z+1|$
        \task $|z-(2+i)|\geq|z+(2+i)|$
    \end{tasks}
\end{exercise}


% ==================================================
% =============== KOMPLEXNÍ JEDNOTKA ================
% ==================================================

\subsection*{Komplexní jednotka}

\begin{exerciseblock}[2]{Rozhodněte, zda je dané číslo komplexní jednotkou}
    \task $\displaystyle \frac12+\frac{\sqrt3}{2}i$
    \task $\displaystyle -\frac{\sqrt2}{2}+\frac{\sqrt2}{2}i$
    \task $\displaystyle \frac{\sqrt6}{5}-\frac{\sqrt{19}}5i$
    \task $\displaystyle \frac{\sqrt6+2\sqrt3}{6}-i\frac{2\sqrt3-\sqrt6}{6}$
    \task $\cos\alpha+i\sin\alpha$
    \task $2(\cos\alpha+i\sin\alpha)$
    \task $\displaystyle \frac{3+4i}{5}$
    \task $\displaystyle \frac{1+i}{\sqrt2}$
    \task $\displaystyle \frac{5-12i}{13}$
    \task $\displaystyle \frac{2+i}{\sqrt5}$
\end{exerciseblock}


\begin{exercise}{Určete všechna reálná čísla $x$, pro která je dané číslo komplexní jednotkou.}
    \begin{tasks}(2)
        \task $\displaystyle x+\frac{\sqrt5}{3}i$
        \task $x+6i$
        \task $\displaystyle x+1+\frac{\sqrt3}{2}i$
        \task $\displaystyle \frac{\sqrt3}{2}+xi+\frac{i}{2}$
        \task $\sin x(\sin x+i)+\cos x(\cos x+1)$
        \task $\displaystyle (1+i)\frac{\sqrt x}{4}+(1-i)\frac{\sqrt5}{4}$
    \end{tasks}
\end{exercise}


% ==================================================
% =============== GONIOMETRICKÝ TVAR ================
% ==================================================

\subsection*{Goniometrický tvar komplexního čísla}

\begin{exerciseblock}[2]{Převeďte do goniometrického tvaru}
    \task $1+i$
    \task $3$
    \task $5i$
    \task $-2+2\sqrt3 i$
    \task $-\sqrt3+i$
    \task $10-10i$
    \task $-7-7i$
    \task $1-\sqrt3 i$
    \task $-4i$
    \task $\sqrt3-3i$
\end{exerciseblock}


\begin{exercise}{Převeďte do goniometrického tvaru. Absolutní hodnotu určete přesně, argument lze v netabulkových případech určit kalkulačkou.}
    \begin{tasks}(3)
        \task $6+3i$
        \task $-1{,}4+5{,}6i$
        \task $9-4i$
        \task $7+11i$
        \task $-8-3i$
        \task $2-9i$
    \end{tasks}
\end{exercise}


\begin{exerciseblock}[2]{Převeďte do algebraického tvaru}
    \task $\displaystyle 4\left(\cos\frac{7\pi}{6}+i\sin\frac{7\pi}{6}\right)$
    \task $\displaystyle \cos\frac{\pi}{2}+i\sin\frac{\pi}{2}$
    \task $5(\cos11\pi+i\sin11\pi)$
    \task $\displaystyle \sqrt2\left(\cos\frac{5\pi}{4}+i\sin\frac{5\pi}{4}\right)$
    \task $\displaystyle 6\left(\cos\frac{\pi}{3}+i\sin\frac{\pi}{3}\right)$
    \task $\displaystyle 3\left(\cos\frac{5\pi}{6}+i\sin\frac{5\pi}{6}\right)$
    \task $\displaystyle 2\left(\cos\frac{3\pi}{2}+i\sin\frac{3\pi}{2}\right)$
    \task $\displaystyle 8\left(\cos\frac{7\pi}{4}+i\sin\frac{7\pi}{4}\right)$
    \task $\displaystyle 10(\cos2\pi+i\sin2\pi)$
    \task $\displaystyle 4\left(\cos\frac{2\pi}{3}-i\sin\frac{\pi}{3}\right)$
\end{exerciseblock}


\begin{exerciseblock}[2]{Upravte a výsledek zapište v goniometrickém tvaru}
    \task $\displaystyle \frac{-1+2i}{1+3i}$
    \task $\displaystyle \frac{1+i}{1-i}$
    \task $\displaystyle \frac{i-3}{2+i}$
    \task $\displaystyle \frac{i-2}{4i-8}$
    \task $(1+i)(1+\sqrt3 i)$
    \task $\displaystyle \frac{2-2i}{\sqrt3+i}$
    \task $(\sqrt3+i)^3$
    \task $\displaystyle \frac{(1+i)^4}{1-i}$
\end{exerciseblock}


% ==================================================
% ===== OPERACE V GONIOMETRICKÉM TVARU =============
% ==================================================

\subsection*{Počítání v goniometrickém tvaru}

\begin{exerciseblock}[1]{Vypočítejte součin a podíl čísel $z_1,z_2$ a výsledek zapište v goniometrickém i algebraickém tvaru}
    \task $z_1=2(\cos105^\circ+i\sin105^\circ)$, $z_2=4(\cos225^\circ+i\sin225^\circ)$
    \task $\displaystyle z_1=2\left(\cos\frac{5\pi}{3}+i\sin\frac{5\pi}{3}\right)$,
    $\displaystyle z_2=4\left(\cos\frac{\pi}{6}+i\sin\frac{\pi}{6}\right)$
    \task $\displaystyle z_1=3\left(\cos\frac{\pi}{4}+i\sin\frac{\pi}{4}\right)$,
    $\displaystyle z_2=2\left(\cos\frac{2\pi}{3}+i\sin\frac{2\pi}{3}\right)$
    \task $\displaystyle z_1=5\left(\cos\frac{7\pi}{6}+i\sin\frac{7\pi}{6}\right)$,
    $\displaystyle z_2=\cos\frac{\pi}{3}+i\sin\frac{\pi}{3}$
\end{exerciseblock}


\begin{exercise}{Je dáno komplexní číslo
\[
    z=6-2i.
\]
Určete komplexní čísla $z_1,z_2,z_3,z_4$ tak, aby platilo:}
    \begin{tasks}(1)
        \task $z_1$ má stejnou absolutní hodnotu a dvojnásobný argument než $z$,
        \task $z_2$ má poloviční absolutní hodnotu a argument o $\frac{5\pi}{4}$ větší než $z$,
        \task $z_3$ má dvojnásobnou absolutní hodnotu a opačný argument než $z$,
        \task $z_4$ má stejnou absolutní hodnotu a argument o $\pi$ větší než $z$.
    \end{tasks}

    Výsledky určete početně a znázorněte v Gaussově rovině.
\end{exercise}


% ==================================================
% ================== UMOCŇOVÁNÍ =====================
% ==================================================

\subsection*{Umocňování komplexních čísel}

\begin{exerciseblock}[2]{Užitím Moivreovy věty umocněte a výsledek převeďte do algebraického tvaru}
    \task $\displaystyle \left(\cos\frac{\pi}{18}+i\sin\frac{\pi}{18}\right)^6$
    \task $\displaystyle \left(\cos\frac{3\pi}{32}+i\sin\frac{3\pi}{32}\right)^8$
    \task $(1+i)^6$
    \task $(1-\sqrt3 i)^5$
    \task $(-2\sqrt3-2i)^{12}$
    \task $(5\sqrt3-5i)^7$
    \task $(\sqrt3+i)^{10}$
    \task $(-1+i)^{16}$
    \task $(2+2i)^{20}$
    \task $(1-\sqrt3 i)^{24}$
\end{exerciseblock}


\begin{exercise}{Určete všechna $n\in\mathbb{N}$, pro která platí}
    \begin{tasks}(2)
        \task $\left(\frac{\sqrt3}{2}-\frac12i\right)^n=
               \frac{\sqrt3}{2}-\frac12i$
        \task $(1+i)^n\in\mathbb{R}$
        \task $(1+i)^n\in i\mathbb{R}$
        \task $\left(\cos\frac{\pi}{6}+i\sin\frac{\pi}{6}\right)^n=1$
        \task $i^n=1$
        \task $(-i)^n=-1$
    \end{tasks}
\end{exercise}


\begin{exercise}{Komplexní číslo
\[
    z=\cos\frac{\pi}{4}+i\sin\frac{\pi}{4}
\]
leží na jednotkové kružnici.}
    \begin{tasks}(2)
        \task Určete $z^8$.
        \task Určete $z^{11}$.
        \task Určete $z^{17}$.
        \task Určete $z^{38}$.
        \task Určete $z^{101}$.
        \task Rozhodněte, kolik různých hodnot může nabývat $z^n$ pro $n\in\mathbb{N}$.
    \end{tasks}

    Všechny různé hodnoty zakreslete do jedné Gaussovy roviny a vysvětlete
    geometrický význam umocňování.
\end{exercise}


% ==================================================
% ================= ODMOCŇOVÁNÍ =====================
% ==================================================

\subsection*{Odmocňování komplexních čísel}

\begin{exerciseblock}[2]{Určete všechny druhé komplexní odmocniny}
    \task z čísla $4$
    \task z čísla $-4$
    \task z čísla $i$
    \task z čísla $-i$
    \task z čísla $3+4i$
    \task z čísla $-3+4i$
    \task z čísla $1+i$
    \task z čísla $-8+6i$
\end{exerciseblock}


\begin{exerciseblock}[2]{Určete všechny uvedené odmocniny}
    \task třetí odmocniny z čísla $8$
    \task třetí odmocniny z čísla $-8$
    \task čtvrté odmocniny z čísla $1$
    \task čtvrté odmocniny z čísla $i$
    \task čtvrté odmocniny z čísla $1-i$
    \task páté odmocniny z čísla $32$
    \task šesté odmocniny z čísla $-64$
    \task osmé odmocniny z čísla $1$
\end{exerciseblock}


\begin{exercise}{Určete všechny kořeny a znázorněte je v Gaussově rovině.}
    \begin{tasks}(2)
        \task $z^3=1$
        \task $z^4=1$
        \task $z^5=1$
        \task $z^6=1$
        \task $z^3=-1$
        \task $z^4=-1$
    \end{tasks}

    U každé rovnice popište geometrický útvar, jehož vrcholy jednotlivé kořeny tvoří.
\end{exercise}


% ==================================================
% ============= ROVNICE V MNOŽINĚ C =================
% ==================================================

\subsection*{Rovnice v množině komplexních čísel}

\begin{exerciseblock}[2]{Řešte v množině $\mathbb{C}$}
    \task $z=3i(z-i)-5z$
    \task $\displaystyle \frac{z}{1-i}+\frac{z+2}{i}=\frac{5}{2i-1}$
    \task $(z+i)(z-3i)=z(z-i)$
    \task $\displaystyle i-\left(\frac{1+i}{1-i}\right)^2=1-\frac1z$
    \task $2z+3\overline z=5+i$
    \task $(2-i)\overline z-13=2(6{,}5i-z)$
    \task $z\overline z-z=6-2i$
    \task $\overline z(z-4)-1=8i$
    \task $z+\overline z=6,\quad z-\overline z=4i$
    \task $|z|^2-2z=3-4i$
\end{exerciseblock}


\begin{exercise}{Řešte v množině $\mathbb{C}$.}
    \begin{tasks}(2)
        \task $|z|=1+2i+z$
        \task $|z+i|=2z+i$
        \task $|z+1|-4i=z+3$
        \task $|z+2-i|=5(z+3i)$
        \task $z+\overline z=4$
        \task $z-\overline z=6i$
        \task $|z|=2$ a $\operatorname{Re}(z)=1$
        \task $|z|=5$ a $\operatorname{Im}(z)=-4$
    \end{tasks}
\end{exercise}


\begin{exerciseblock}[2]{Vyřešte soustavu v množině komplexních čísel}
    \task $\begin{cases}2z+w=8i,\\z-w=6+i,\end{cases}$
    \task $\begin{cases}(1+i)z-3w=-7-6i,\\2z+(2+i)w=6+5i,\end{cases}$
    \task $\begin{cases}z+w=3+i,\\z-w=1-5i,\end{cases}$
    \task $\begin{cases}iz+w=2,\\z-iw=4i.\end{cases}$
\end{exerciseblock}


% ==================================================
% ============== KVADRATICKÉ ROVNICE ================
% ==================================================

\subsection*{Kvadratické rovnice v množině komplexních čísel}

\begin{exerciseblock}[2]{Řešte v množině $\mathbb{C}$}
    \task $x^2-5x+5=0$
    \task $5x^2-2x+1=0$
    \task $3x^2-2x+1=0$
    \task $x^2-4ix-8=0$
    \task $x^2-6ix-9=0$
    \task $x^2+4x+5=0$
    \task $x^2-2x+10=0$
    \task $9x^2-6x+5=0$
    \task $x^2+(2-i)x+2=0$
    \task $x^2-(3+i)x+4i=0$
\end{exerciseblock}


\begin{exerciseblock}[2]{Řešte v množině $\mathbb{C}$}
    \task $x^2+x(2-i)+3-i=0$
    \task $(7+i)x^2-5ix-1=0$
    \task $x(3-x)=3-i$
    \task $x^2-20=ix(2i-x)$
    \task $ix^2-3x+4i=0$
    \task $x^2+(i-3)x+2-2i=0$
    \task $(1+i)x^2+(2-i)x+1=0$
    \task $(2-3i)x^2+ix-4=0$
\end{exerciseblock}


\begin{exercise}{Sestavte kvadratickou rovnici s reálnými koeficienty, znáte-li jeden její kořen.}
    \begin{tasks}(2)
        \task $x_1=3+i$
        \task $x_1=5i$
        \task $\displaystyle x_1=\cos\frac{5\pi}{3}+i\sin\frac{5\pi}{3}$
        \task $\displaystyle x_1=2\left(\cos\frac{11\pi}{6}+i\sin\frac{11\pi}{6}\right)$
        \task $x_1=-2+3i$
        \task $x_1=1-\sqrt2 i$
    \end{tasks}
\end{exercise}


\begin{exercise}{Určete druhý kořen a neznámý koeficient.}
    \begin{tasks}(1)
        \task Rovnice $x^2+ix+q=0$ má kořen $x_1=2-i$.
        \task Rovnice $x^2+px+21=0$ má kořen $x_1=-3+2i\sqrt3$.
        \task Rovnice $x^2+2(3-2i)x+k=0$ má dvojnásobný kořen.
        \task Rovnice $x^2+px+13=0$ má jeden kořen $2+3i$.
    \end{tasks}
\end{exercise}


% ==================================================
% ================ BINOMICKÉ ROVNICE ================
% ==================================================

\subsection*{Binomické rovnice}

\begin{exerciseblock}[2]{Řešte v množině $\mathbb{C}$, kořeny zapište v goniometrickém i algebraickém tvaru a znázorněte je v Gaussově rovině}
    \task $z^3-27=0$
    \task $z^4+16=0$
    \task $z^6-1=0$
    \task $z^3-64i=0$
    \task $z^4-16=0$
    \task $z^5+32=0$
    \task $z^3+8i=0$
    \task $z^6+64=0$
\end{exerciseblock}


\begin{exerciseblock}[2]{Řešte v množině $\mathbb{C}$}
    \task $z^3-1-i=0$
    \task $z^6-1+i\sqrt3=0$
    \task $(iz)^4+\sqrt3-i=0$
    \task $(2z)^5-16=16i\sqrt3$
    \task $z^4=8+8\sqrt3 i$
    \task $z^5=-32i$
    \task $z^8=1$
    \task $z^{10}=-1$
\end{exerciseblock}


% ==================================================
% =================== BONUS ==========================
% ==================================================

\subsection*{Kombinované a náročnější úlohy}

\begin{exercise}{Dokažte, že pro každé reálné číslo $p$ je číslo
\[
    z=\frac{i}{p-3i}+\frac{i}{p+3i}
\]
ryze imaginární. Poté určete všechna $p\in\mathbb{R}$, pro která platí
\[
    z=\frac35i.
\]}
\end{exercise}


\begin{exercise}{Dokažte, že podíl
\[
    \frac{m-3i}{3+mi}
\]
nezávisí na hodnotě $m\in\mathbb{R}$, pokud je výraz definován.}
\end{exercise}


\begin{exercise}{Určete všechna komplexní čísla $z$, pro která současně platí
\[
    |z|=2
\]
a
\[
    |z-2|=2.
\]
Řešení určete početně i geometricky.}
\end{exercise}


\begin{exercise}{Je dáno komplexní číslo
\[
    z=\frac{\sqrt3}{2}+\frac12i.
\]
Bez přímého roznásobování určete
\[
    z^{10},\qquad z^{25},\qquad z^{100},\qquad z^{2026}.
\]
Vysvětlete výsledky pomocí rotace v Gaussově rovině.}
\end{exercise}


\begin{exercise}{Najděte všechna komplexní čísla $z$, pro která platí
\[
    z^4=\overline z^4
\]
a současně
\[
    |z|=2.
\]
Řešení znázorněte v Gaussově rovině.}
\end{exercise}


\begin{exercise}{Určete všechna komplexní čísla $z$, která splňují
\[
    z^3=\overline z.
\]
}
\end{exercise}


\begin{exercise}{Určete všechny kořeny rovnice
\[
    z^8=1.
\]
Poté určete součet všech kořenů, jejich součin a součet jejich druhých mocnin.}
\end{exercise}


\begin{exercise}{Komplexní číslo $z$ splňuje
\[
    |z|=1
\]
a
\[
    z+\overline z=1.
\]
Určete všechny možné hodnoty $z$ a poté vypočítejte $z^{2026}$.}
\end{exercise}


% ==================================================
% ============= DALŠÍ ÚLOHY K PROCVIČENÍ ===========
% ==================================================

\subsection*{Normální úlohy}


% --------------------------------------------------
% REÁLNÉ A NEREÁLNÉ HODNOTY
% --------------------------------------------------

\begin{exerciseblock}[2]{Vypočítejte hodnotu výrazu a rozhodněte, zda výsledné číslo patří do množiny $\mathbb{R}$}
    \task $(2+i)(2-i)$
    \task $\pi\cdot i^{16}$
    \task $\displaystyle \left(\cos\frac{\pi}{2}+i\sin\frac{\pi}{2}\right)^2$
    \task $\displaystyle \left(\frac{1}{i-1}\right)^2$
    \task $\displaystyle i+\frac{1}{i}$
    \task $(1+i)^4$
    \task $(2-i)^2$
    \task $\displaystyle \left(\frac{1+i}{1-i}\right)^2$
    \task $(3+2i)+(3-2i)$
    \task $\displaystyle \frac{1+i}{2-i}$
\end{exerciseblock}


% --------------------------------------------------
% ABSOLUTNÍ HODNOTA A NEROVNOST
% --------------------------------------------------

\begin{exercise}{Pro každé z následujících komplexních čísel rozhodněte, zda platí
\[
    |z+3i|\leq4.
\]
Své rozhodnutí zdůvodněte výpočtem.}

    \begin{tasks}(2)
        \task $z=-7i$
        \task $z=-4$
        \task $z=3-5i$
        \task $z=4-3i$
        \task $z=-4-3i$
        \task $z=1+i$
        \task $z=-3i$
        \task $z=2-6i$
        \task $z=3+i$
        \task $z=-2-6i$
    \end{tasks}
\end{exercise}


\begin{exerciseblock}[2]{Vypočítejte}
    \task $\displaystyle \frac{|3i|+7}{|3-4i|}$
    \task $\displaystyle \frac{|1+i|}{|2-i|}$
    \task $\displaystyle \frac{|3+4i|}{|-2i|}$
    \task $\displaystyle \frac{|2-2i|}{|1+i|}$
    \task $\displaystyle \left|\frac{2+3i}{1-i}\right|$
    \task $\displaystyle \left|\frac{4-i}{2+2i}\right|$
    \task $\displaystyle \frac{|(1+i)(3-4i)|}{|2-i|}$
    \task $\displaystyle \frac{|(2+i)^2|}{|1-2i|}$
    \task $\displaystyle \left|\frac{(1-i)^3}{2+i}\right|$
    \task $\displaystyle \frac{|z_1z_2|}{|z_1|}$, je-li $|z_2|=5$
\end{exerciseblock}


% --------------------------------------------------
% ALGEBRAICKÝ TVAR
% --------------------------------------------------

\begin{exerciseblock}[2]{Vyjádřete komplexní číslo v algebraickém tvaru}
    \task $\displaystyle z=\frac{5i}{2+i}$
    \task $\displaystyle z=\frac{3}{1-2i}$
    \task $\displaystyle z=\frac{2+i}{1-i}$
    \task $\displaystyle z=\frac{1}{2i}$
    \task $\displaystyle z=\frac{3-4i}{1+i}$
    \task $\displaystyle z=\frac{2}{3+i}$
    \task $\displaystyle z=\frac{(1+i)^2}{2-i}$
    \task $\displaystyle z=\frac{4i}{1+2i}$
    \task $\displaystyle z=\frac{5-3i}{2+i}$
    \task $\displaystyle z=\frac{1+i}{1-i}+\frac{1-i}{1+i}$
\end{exerciseblock}


% --------------------------------------------------
% KOMPLEXNĚ SDRUŽENÁ ČÍSLA
% --------------------------------------------------

\begin{exerciseblock}[2]{Řešte rovnici s neznámou $z\in\mathbb{C}$}
    \task $z+3i=2\overline{z}-4$
    \task $z-2\overline{z}=1+3i$
    \task $2z+\overline{z}=9-i$
    \task $3z-2\overline{z}=-4+5i$
    \task $2\overline{z}-z=3-6i$
    \task $4z-\overline{z}=6+15i$
    \task $z+\overline{z}=8,\quad z-\overline{z}=4i$
    \task $2z+\overline{z}=3,\quad z-\overline{z}=-2i$
\end{exerciseblock}


% --------------------------------------------------
% GAUSSOVA ROVINA -- MNOŽINY BODŮ
% --------------------------------------------------

\begin{exercise}{V Gaussově rovině znázorněte množinu všech komplexních čísel $z$, pro která platí}
    \begin{tasks}(2)
        \task $|z|=3$
        \task $z+\overline{z}=-2$
        \task $|z-(2-i)|=2$
        \task $|z+1|\leq|z-i|$
        \task $|z-(1+i)|=|z-(3-i)|$
        \task $|z-2i|<3$
        \task $z-\overline{z}=4i$
        \task $|z-1|=|z+3|$
        \task $|z-(1+i)|\leq2$
        \task $2<|z+2i|\leq4$
    \end{tasks}
\end{exercise}


\begin{exercise}{U každé z předchozích množin určete, zda se jedná o bod, přímku, kružnici, kruh, polorovinu nebo jiný geometrický útvar. Je-li to možné, určete také jeho střed, poloměr nebo rovnici příslušné přímky.}
\end{exercise}


% --------------------------------------------------
% OPERACE A GAUSSOVA ROVINA
% --------------------------------------------------

\begin{exercise}{Jsou dána komplexní čísla
\[
    z_1=2-2i,
    \qquad
    z_2=-2.
\]
Nejprve zobrazte jejich obrazy v Gaussově rovině a následně určete a zakreslete obrazy čísel}
    \begin{tasks}(4)
        \task $z_1+z_2$
        \task $z_1-z_2$
        \task $z_1z_2$
        \task $\displaystyle \frac{z_1}{z_2}$
    \end{tasks}
\end{exercise}


\begin{exercise}{Pro každou dvojici komplexních čísel vypočítejte podíl
\[
    z=\frac{z_1}{z_2}
\]
a čísla $z_1$, $z_2$ a $z$ znázorněte v jedné Gaussově rovině.}

    \begin{tasks}(2)
        \task $z_1=2-2i,\quad z_2=-2$
        \task $z_1=1+3i,\quad z_2=i$
        \task $z_1=-2+2i,\quad z_2=1-i$
        \task $z_1=3+i,\quad z_2=1+i$
        \task $z_1=-4i,\quad z_2=2i$
        \task $z_1=2+2i,\quad z_2=-1+i$
    \end{tasks}
\end{exercise}


% --------------------------------------------------
% ODMOCNINY A JEDNODUCHÉ ROVNICE
% --------------------------------------------------

\begin{exercise}{Řešte v množině $\mathbb{C}$ a všechny kořeny znázorněte v Gaussově rovině.}
    \begin{tasks}(2)
        \task $z^2=-i$
        \task $z^2=i$
        \task $z^2=-4$
        \task $z^2=4i$
        \task $z^3=8$
        \task $z^3=-8$
        \task $z^4=1$
        \task $z^4=-1$
    \end{tasks}
\end{exercise}


% --------------------------------------------------
% GONIOMETRICKÝ TVAR A MOCNINY
% --------------------------------------------------

\begin{exerciseblock}[2]{Vypočítejte a výsledek zapište v algebraickém tvaru}
    \task $\displaystyle
        z=\frac12\left(
            \cos\frac{4\pi}{9}
            +i\sin\frac{4\pi}{9}
        \right),
        \qquad
        \frac1{z^3}$
    \task $\displaystyle
        z=2\left(
            \cos\frac{\pi}{6}
            +i\sin\frac{\pi}{6}
        \right),
        \qquad
        z^3$
    \task $\displaystyle
        z=\sqrt2\left(
            \cos\frac{3\pi}{4}
            +i\sin\frac{3\pi}{4}
        \right),
        \qquad
        z^4$
    \task $\displaystyle
        z=\frac13\left(
            \cos\frac{5\pi}{6}
            +i\sin\frac{5\pi}{6}
        \right),
        \qquad
        \frac1{z^2}$
    \task $\displaystyle
        z=2\left(
            \cos\frac{2\pi}{3}
            +i\sin\frac{2\pi}{3}
        \right),
        \qquad
        z^6$
    \task $\displaystyle
        z=3\left(
            \cos\frac{7\pi}{6}
            +i\sin\frac{7\pi}{6}
        \right),
        \qquad
        \frac1{z^2}$
\end{exerciseblock}


\begin{exercise}{Jsou dána komplexní čísla
\[
    z_1=\frac{1}{4i}
\]
a
\[
    z_2=
    2\left(
        \cos\frac{11\pi}{6}
        -
        i\sin\frac{11\pi}{6}
    \right).
\]

Bez převodu obou čísel do algebraického tvaru určete

\[
    \left|\frac{z_2}{z_1}\right|.
\]
}
\end{exercise}


\begin{exerciseblock}[2]{Vypočítejte pouze pomocí vlastností absolutní hodnoty komplexních čísel}
    \task $\displaystyle \left|\frac{3+4i}{1-i}\right|$
    \task $\displaystyle \left|\frac{(1+i)^5}{2-i}\right|$
    \task $\displaystyle \left|\frac{z_1^4}{z_2^3}\right|$, je-li $|z_1|=2$ a $|z_2|=4$
    \task $\displaystyle |z_1^3z_2^2|$, je-li $|z_1|=3$ a $|z_2|=\frac12$
    \task $\displaystyle \left|\frac{\overline z}{z}\right|$, kde $z\neq0$
    \task $\displaystyle |z\overline z|$, je-li $|z|=5$
\end{exercise}



% ==================================================
% ==================== BONUS =========================
% ==================================================

\subsection*{BONUS úlohy}


% --------------------------------------------------
% GEOMETRICKÁ INTERPRETACE
% --------------------------------------------------

\begin{exercise}{BONUS: Pro každou z následujících podmínek určete geometrický význam množiny komplexních čísel $z$, odvoďte její rovnici v souřadnicích a množinu zakreslete v Gaussově rovině.}

    \begin{tasks}(1)
        \task
        \[
            |z-(1+i)|=|z-(3-i)|
        \]

        \task
        \[
            |z+1|\leq|z-i|
        \]

        \task
        \[
            |z-(2-i)|=2
        \]

        \task
        \[
            |z-(2+i)|<|z-(-2+i)|
        \]

        \task
        \[
            |z-(1-2i)|=|z+(3+i)|
        \]

        \task
        \[
            |z-2|+|z+2|=6
        \]
    \end{tasks}
\end{exercise}


\begin{exercise}{BONUS: Vytvořte v Gaussově rovině vlastní příklad}
    \begin{tasks}(1)
        \task přímky zadané pomocí rovnosti vzdáleností od dvou komplexních čísel,
        \task poloroviny zadané pomocí nerovnosti dvou vzdáleností,
        \task kružnice se středem mimo obě souřadnicové osy,
        \task mezikruží zadaného pomocí absolutní hodnoty.
    \end{tasks}

    Ke každému obrázku napište odpovídající podmínku s komplexním číslem $z$.
\end{exercise}


% --------------------------------------------------
% ČTVRTÉ ODMOCNINY
% --------------------------------------------------

\begin{exercise}{BONUS: V oboru $\mathbb{C}$ řešte rovnici
\[
    z^4=(2i)^2.
\]

Všechna řešení nejprve určete v goniometrickém tvaru, poté je převeďte
do algebraického tvaru a nakonec jejich obrazy znázorněte v Gaussově rovině.
}
\end{exercise}


\begin{exercise}{BONUS: Řešte v množině $\mathbb{C}$ a všechny kořeny zakreslete do jedné Gaussovy roviny.}
    \begin{tasks}(2)
        \task $z^4=-4$
        \task $z^6=-64$
        \task $z^5=32i$
        \task $z^8=16$
        \task $z^6=8i$
        \task $z^{10}=-1$
    \end{tasks}
\end{exercise}


% --------------------------------------------------
% NEJMENŠÍ EXPONENT
% --------------------------------------------------

\begin{exercise}{BONUS: Je dáno komplexní číslo
\[
    z=\sqrt3-i.
\]

Najděte nejmenší kladné celé číslo $m$, pro které je číslo

\[
    z^m
\]

celé, tedy $z^m\in\mathbb{Z}$, a následně hodnotu $z^m$ vypočítejte.
}
\end{exercise}


\begin{exerciseblock}[2]{BONUS: Určete nejmenší kladné celé číslo $m$, pro které platí $z^m\in\mathbb{Z}$, a vypočítejte $z^m$}
    \task $z=1+i$
    \task $z=1+\sqrt3 i$
    \task $z=\sqrt3+i$
    \task $z=-1+i$
    \task $z=\sqrt3-i$
    \task $z=-\sqrt3+i$
    \task $z=1-\sqrt3 i$
    \task $z=-1-\sqrt3 i$
\end{exerciseblock}


% --------------------------------------------------
% MOCNINY A PERIODICITA
% --------------------------------------------------

\begin{exercise}{BONUS: Je dáno komplexní číslo
\[
    z=\sqrt3-i.
\]

Bez postupného roznásobování určete hodnoty

\[
    z^6,\qquad
    z^{12},\qquad
    z^{17},\qquad
    z^{25},\qquad
    z^{100}.
\]

U každé mocniny určete její absolutní hodnotu a hlavní argument a vysvětlete,
jak se při umocňování mění poloha výsledného čísla v Gaussově rovině.
}
\end{exercise}


\begin{exercise}{BONUS: Je dána komplexní jednotka
\[
    z=
    \cos\frac{\pi}{6}
    +
    i\sin\frac{\pi}{6}.
\]

Určete

\[
    z^8,\qquad
    z^{11},\qquad
    z^{17},\qquad
    z^{38},\qquad
    z^{100},\qquad
    z^{2026}.
\]

Poté určete počet všech různých hodnot, kterých může nabývat $z^n$ pro
$n\in\mathbb{N}$, a všechny tyto hodnoty zakreslete do jedné Gaussovy roviny.
}
\end{exercise}


% --------------------------------------------------
% KOMPLEXNĚ SDRUŽENÉ ČÍSLO
% --------------------------------------------------

\begin{exerciseblock}[2]{BONUS: Řešte v množině $\mathbb{C}$}
    \task $z+3i=2\overline z-4$
    \task $(1+i)z+(1-i)\overline z=6$
    \task $z^2+\overline z^2=8$
    \task $z\overline z-z-\overline z=3$
    \task $z+\overline z=|z|^2$
    \task $z^2=\overline z$
\end{exerciseblock}


% --------------------------------------------------
% KOMBINACE MODULU A GEOMETRIE
% --------------------------------------------------

\begin{exercise}{BONUS: Určete všechna komplexní čísla $z$, která současně splňují
\[
    |z+3i|\leq4
\]
a
\[
    |z-1|=|z+3|.
\]

Množinu všech řešení znázorněte v Gaussově rovině.
}
\end{exercise}


\begin{exercise}{BONUS: Určete všechna komplexní čísla $z$, která splňují
\[
    |z-(1+i)|=2
\]
a současně
\[
    z+\overline z=2.
\]

Řešení určete početně i geometricky.
}
\end{exercise}


\begin{exercise}{BONUS: Určete všechna komplexní čísla $z$, pro která platí
\[
    |z|=|z-4|
\]
a současně
\[
    |z-2i|=\sqrt5.
\]

Všechna řešení zakreslete v Gaussově rovině.
}
\end{exercise}


% --------------------------------------------------
% KOMBINOVANÉ ÚLOHY
% --------------------------------------------------

\begin{exercise}{BONUS: Je dáno komplexní číslo
\[
    z=
    \frac12
    \left(
        \cos\frac{4\pi}{9}
        +
        i\sin\frac{4\pi}{9}
    \right).
\]

Bez převodu čísla $z$ do algebraického tvaru vypočítejte

\[
    \frac{1}{z^3}
\]

a výsledek následně zapište v algebraickém tvaru.
}
\end{exercise}


\begin{exercise}{BONUS: Jsou dána komplexní čísla
\[
    z_1=\frac{1}{4i}
\]

a

\[
    z_2=
    2\left(
        \cos\frac{11\pi}{6}
        -
        i\sin\frac{11\pi}{6}
    \right).
\]

Určete bez úplného převodu obou čísel do algebraického tvaru

\[
    \left|\frac{z_2}{z_1}\right|,
\]

hlavní argument čísla

\[
    \frac{z_2}{z_1}
\]

a nakonec tento podíl zapište v algebraickém tvaru.
}
\end{exercise}


\begin{exercise}{BONUS: Je dáno nenulové komplexní číslo $z$, pro které platí
\[
    |z|=2
\]

a

\[
    \operatorname{Arg}z=\frac{\pi}{6}.
\]

Určete nejmenší kladné celé číslo $n$, pro které je $z^n$ reálné, a nejmenší
kladné celé číslo $m$, pro které je $z^m$ kladné reálné číslo. Obě příslušné
mocniny vypočítejte.
}
\end{exercise}
